% !TeX root = ../教材题库_全部条目_转写工作稿.tex
% Source: XJ-B1-C04-TH-P109-Y0302
% 第4章 / 4.2 指数函数 / 教材 p.109、110 / PDF 第116、117页
% OCR未校对：公式与文字均可能有误；未完成数学TeX转写；禁止直接发给学生
% 原题及相邻原解截图（核对用）：../题目原图/XJ-B1-C04-TH-P109-Y0302-01.webp, ../题目原图/XJ-B1-C04-TH-P109-Y0302-02.webp
% 本文件只在转写完成后作为题干使用；OCR 条目尚不是可靠题干。
\begin{bankitem}{p.109 思考（含紧邻说明）}{第4章 / 4.2 指数函数 / 教材 p.109、110 / PDF 第116、117页}{XJ-B1-C04-TH-P109-Y0302}{OCR未校对：公式与文字均可能有误；未完成数学TeX转写；禁止直接发给学生}
\begin{ocrdraft}
[教材p.109 / PDF第116页]
\par
？思考
\par
下表是指数函数3，一2'和一2
\par
在自变量x分别取一些数值时所对
\par
应的函数值．
\par
一100
\par
0
\par
5
\par
10
\par
100
\par
1
\par
1
\par
1
\par
1
\par
32
\par
21
\par
21
\par
32
\par
1
\par
1
\par
1
\par
21
\par
21
\par
32
\par
1
\par
2
\par
（1）观察上表，你能猜测出函数一2'和3一
\par
一）的值域与单调性吗？
\par
（2）上述猜测是否可以看作指数函数3，一的性质？
\par
（1）由上表可以猜测：3，一2x的函数值均大于0，且随着自变量工的增大而增
\par
大；对于3一），其自变量越小，函数值越大，自变量越大，函数值越小，
\par
但总大于0．同时，两函数的图象均过点（0，1）．
\par
（2）上述猜测可以分别看作指数函数3，一>1)或3，一（0<a<1)的性质．
\par
理由如下：
\par
由于任意非零实数的零次幂为1，即一1，因此任意指数函数的图象都经过点
\par
又由于正数的任意次幂仍为正，因而任意指数函数的图象都在工轴上方（即函数
\par
值均大于0）．
\par
另外，由前面所学的对实数成立的幂运算的基本不等式可得：
\par
对任意实数。厂，s，若磊>1，7．>、，则一
\par
a冖、>1，即>．
\par
对任意实数厂，s，若0<磊<1，7．>、，贝刂一一a一'<1，即<．
\par
又指数函数的定义域为R，
\par
于是，当磊>1时，指数函数3，一a'在（一。。，十。。）上单调递增；
\par
当0<<1时，指数函数3-a'在（一。。，十亠）上单调递减．
\par
我们还可以证明：当a>1时，令a一1十方，则禳'一（1十/D">劢0N十）（D恒成
\par
1
\par
\ensuremath{\cdot}这表明，只要足够大，a'可以大于任意给定的正数，而a-
\par
一则可以小于
\par
(D利用选择性必修第一册中的相关知识可以证明此结论．
\par
[教材p.110 / PDF第117页]
\par
任意给定的正数．综上可得，指数函数的值域为（0，十。。）．
\par
至此，我们对指数函数的性质有了基本的理性认识，下面从直观的角度探讨指
\par
数函数的性质．
\end{ocrdraft}
\end{bankitem}
